\documentclass[11pt]{amsart}
\usepackage[margin=1in]{geometry}
\usepackage{amsmath,amssymb,amsthm,mathtools,booktabs}
\usepackage{enumitem}
\usepackage{hyperref}
\hypersetup{colorlinks=true,linkcolor=blue,citecolor=blue,urlcolor=blue}
\usepackage[final]{microtype}
\setlength{\emergencystretch}{3em}

\newcommand{\leanrepo}{https://github.com/HQIV/hqiv-lean/blob/main}
% Lean module / identifier helpers: visible text is rendered with \nolinkurl
% so that slashes, underscores, and dots become legal break points.
\newcommand{\leanfile}[1]{\href{\leanrepo/#1}{\nolinkurl{#1}}}
\newcommand{\leanid}[1]{\nolinkurl{#1}}

\title[Finite-Mode Kirchhoff from HQIV]
{Kirchhoff's Law of Thermal Emission with Built-In UV/IR Cutoffs
from HQIV's Discrete Null Lattice}

\author{Steven Ettinger Jr}
\address{Independent Researcher}
\email{steven@disregardfiat.tech}
\date{Preprint v1 --- \today}

\newtheorem{definition}{Definition}[section]
\newtheorem{proposition}[definition]{Proposition}
\newtheorem{theorem}[definition]{Theorem}
\newtheorem{lemma}[definition]{Lemma}
\newtheorem{remark}[definition]{Remark}
\newtheorem{corollary}[definition]{Corollary}

\begin{document}
\begin{abstract}
\noindent\textbf{Position in the publication chain.}
This is the third paper in the Tier-1 foundation-extension trilogy:
the HQIV main paper \cite{ettinger2026hqiv} fixes the framework, the
lightcone paper \cite{hqiv-lightcone-paper} introduces the discrete
null-lattice readout, the closure paper \cite{hqiv-so8-paper}
certifies the $G_2\cup\{\Delta\}\Rightarrow\mathfrak{so}(8)$ Lie
closure, the 3D causal-growth paper \cite{hqiv-3d-growth-paper}
isolates the dimensional hinge and the rational \((\alpha_d,\gamma_d)\)
family, and the octonionic-action paper \cite{hqiv-oct-action-paper}
packages the variational layer that we ride here for the
$\alpha=3/5$ imprint and the patch-ontology reader contract. The
Tier-0 records and the 3D causal-growth companion
\cite{hqiv-3d-growth-paper} carry minted Zenodo DOIs; the
octonionic-action record cited above remains pre-DOI (this manuscript
and the remaining Tier-1 queue co-deposit as the Lean PR builds clear).

\noindent\textbf{Result.}
Kirchhoff's law of thermal emission is recovered from two HQIV
axioms---a discrete null lattice with Planck cell size and
informational monogamy ($\alpha=3/5$)---with a \emph{finite} mode
count at every shell, given in closed form by the stars-and-bars
identity
\[
  \mathcal{N}(m)=(m+2)(m+1),
\qquad
  \sum_{m=0}^{n}\mathcal{N}(m)=\tfrac{1}{3}(n+1)(n+2)(n+3).
\]
The resulting Planck spectrum is a \emph{finite sum} between an
ultraviolet cutoff at $m=0$ (the Planck pole, $\omega_{\rm Pl}=T_{\rm Pl}$)
and an arbitrary infrared cutoff at $m=M$ (the lock-in reference
shell).  No ultraviolet catastrophe arises, no renormalisation is
performed, and no external blackbody calibration is imposed: the
mode budget is \emph{finite by construction}, not by regularisation.
This is fully consistent with the patch-theory reader contract
imported below \cite{hqiv-oct-action-paper}: the discrete patch
layer is the ontology, and what the literature calls ``the
continuum limit'' is admitted only as a calculation-approximation
rewrite for cross-citation, never as a foundation.

\noindent\textbf{Sourcing of the dimensionful inputs.}
The proton mass at the lock-in shell $m=4$ uses the theory's
internal constituent-plus-binding readout, with the observed
value $M_{p}=938.272\,\mathrm{MeV}$ serving as the
\emph{active scale witness} (the \hyperref[lean:proton-lockin]{proton lock-in scale-witness} convention; see
\texttt{AGENTS/ASSUMPTIONS.md} \S7g of the HQIV-LEAN repository for
the witness policy).
The wall-clock age $t_{\rm wall}=51.2\,\mathrm{Gyr}$ is an
\emph{output} of the HQIV-modified CLASS implementation---a
discrete null-lattice background with the ADM lapse
$N(t)=1+\Phi+\varphi t$ and $\varphi$-dependent evolution (the
\hyperref[lean:wall-clock]{HQVM wall-clock compression})---not a
free parameter.  The only externally measured dimensionful input
used downstream of these is FIRAS $T_{\rm CMB}=2.7255\,\mathrm{K}$
\cite{fixsen2009}.

\noindent\textbf{One quantitative prediction under the chosen
witness.}  With the proton lock-in scale-witness fixing the
mass/unit chart, all other dimensionful quantities---CODATA
$1/\alpha$, CMB horizons, cosmic birefringence---become
\emph{predictions or comparison layers}.  We exhibit one such
prediction here:
\begin{itemize}[leftmargin=2em]
\item \textbf{Isotropic cosmic birefringence at present epoch.}
The birefringence calculation uses the HQIV wall-clock age
above together with the lattice mode count $\mathcal{N}(m)$
under the propagation-shell identification (\S\ref{sec:beta}),
giving $\beta_{\rm HQIV}=0.379^\circ$. This agrees with Planck
PR4 EB ($\beta_{\rm obs}=0.342^\circ\pm 0.094^\circ$
\cite{eskilt2022}) within $-0.40\sigma$. Substituting HQIV's
apparent age $13.8\,\mathrm{Gyr}$ for the wall-clock age gives
$0.103^\circ$ ($-2.55\sigma$): the prediction
\emph{discriminates} HQIV's two internal time outputs.
\end{itemize}

\noindent\textbf{Status.}  The derivation chain is mechanised
sorry-free against mathlib in the public slimmed Lean mirror
\cite{hqiv-lean}; reader-friendly anchors with clickable Lean
links live in Appendix~\ref{app:lean-index}. The Friedmann
radiation-era scaling $H\propto T^{2}$ is also recovered, not
imported: it follows from the quadratic asymptotic form of
$\mathcal{N}(m)$ under the propagation-shell identification. All
numerical results in this paper (the $\beta=0.379^\circ$
prediction, the wall-clock-vs-apparent-age discriminator, and the
finite-mode Kirchhoff energy density at FIRAS $T_{\rm CMB}$) are
reproduced by the supplementary script bundle described in
Appendix~\ref{app:scripts}, packaged as \texttt{scripts.zip} for
co-deposit with the manuscript on Zenodo.
\end{abstract}
\maketitle

% Shared reader contract: patch QFT + discrete numerics.
% From papers/*.tex:     % Shared reader contract: patch QFT + discrete numerics.
% From papers/*.tex:     % Shared reader contract: patch QFT + discrete numerics.
% From papers/*.tex:     \input{include/patch_theory_messaging.tex}
% From papers/paper/*.tex: \input{../include/patch_theory_messaging.tex}

\paragraph{Patch QFT and discrete numerics (reader contract).}
HQIV (Horizon-Quantized Informational Vacuum) is formulated as a \emph{patch theory}: quantum fields, actions, and physical readouts are attached to discrete patches---shell indices $m\in\mathbb{N}$, finite spacetime index types (e.g.\ $\mathrm{Fin}\,4$), and horizon-limited charts---rather than to a hypothetical fundamental smooth spacetime obtained as a mesh-refinement or ultraviolet limit.
Accordingly, when this text refers to ``QFT,'' it means \emph{patch QFT}: patching rules, finite measurement ledgers, and variational identities on the discrete spine; it is not the claim that the theory is \emph{defined} by recovering ordinary continuum quantum field theory through such a limit.
The numbers that admit the sharpest statements---mass and coupling ladders, curvature ratios, shell thermodynamics, and rational witnesses in \texttt{HQIV\_LEAN}---are objects of \emph{discrete mathematics} on $\mathbb{N}$ (combinatorial growth, cumulative simplex ledgers) together with the ordered field $\mathbb{R}$ as used in the proof assistant for analysis, bounds, and phenomenological display.
Smooth-manifold or continuum field notation, when it appears, is a \emph{translation layer} for comparison with standard physics literature; it does not supersede the discrete definitions.

\paragraph{A continuum is not required, and in places would destroy these results.}
The discrete patch layer is \emph{not} a numerical approximation to a more fundamental continuum theory awaiting future construction; it is the ontology.  Where the literature treats ``the continuum limit'' as the goal, HQIV treats it as a \emph{calculation approximation} for comparison with standard formulas.  The continuum form is an optional convenience for comparison; the proved discrete statement is what carries the physics.  In several load-bearing places---the curvature imprint $\delta_E(m)$ at shell $m$, the harmonic ladder masses $m_\tau\to m_\mu\to m_e$, the lock-in vev $v=\sqrt{\eta_{\rm paper}\,\Omega_k(m_{\rm lockin};m_{\rm lockin})}$, the proton anchor at $m=\mathrm{referenceM}=4$, the baryon-asymmetry $\eta\propto 6^7\sqrt{3}/(m+1)$, and the rational coefficients $\alpha=3/5$, $\gamma=2/5$, $\alpha_{\rm GUT}=1/42$---taking a literal $m\to\infty$ or sub-Planck continuum limit would \emph{erase} the discrete index that is doing the predictive work.  Continuum forms of these objects are therefore not preferred refinements; they are degenerate, information-losing readouts of the same patch data.

\paragraph{An exception: $\mathfrak{so}(8)$ as a de-facto continuum algebra from a discrete construction.}
Principled exception to the discrete-only stance: the closed Lie algebra $\mathfrak{so}(8) \supseteq G_2\cup\{\Delta\}$ generated from the Fano octonion left-multiplication table and the phase-lift generator $\Delta\in(\mathrm{e}_1,\mathrm{e}_7)$.  Although $\mathfrak{so}(8)$ is a continuous (28-dimensional real) Lie algebra, it is built here entirely from \emph{discrete} ingredients (the seven imaginary octonion units, a finite Fano table, and one $\pi/2$ rotation), and its closure to 28 dimensions is a finite linear-algebra certificate (see the \texttt{Hqiv.SO8Closure} / \texttt{Hqiv.SO8ClosureInterface} stack).  The Lie-group structure is therefore a \emph{de-facto continuum algebra obtained from a discrete construction}: continuous in parameter space, but with no continuum spacetime, no continuum measure, and no UV limit attached.  When this manuscript or its companions write $\exp(t\,\Delta)\in\mathrm{SO}(8)$ or treat $\mathfrak{so}(8)$ rotations as smooth, those are statements internal to a finite-dimensional Lie group whose generators are certified by discrete data; they are \emph{not} a continuum-spacetime commitment and do not require a continuum-QFT scaffolding.

% From papers/paper/*.tex: % Shared reader contract: patch QFT + discrete numerics.
% From papers/*.tex:     \input{include/patch_theory_messaging.tex}
% From papers/paper/*.tex: \input{../include/patch_theory_messaging.tex}

\paragraph{Patch QFT and discrete numerics (reader contract).}
HQIV (Horizon-Quantized Informational Vacuum) is formulated as a \emph{patch theory}: quantum fields, actions, and physical readouts are attached to discrete patches---shell indices $m\in\mathbb{N}$, finite spacetime index types (e.g.\ $\mathrm{Fin}\,4$), and horizon-limited charts---rather than to a hypothetical fundamental smooth spacetime obtained as a mesh-refinement or ultraviolet limit.
Accordingly, when this text refers to ``QFT,'' it means \emph{patch QFT}: patching rules, finite measurement ledgers, and variational identities on the discrete spine; it is not the claim that the theory is \emph{defined} by recovering ordinary continuum quantum field theory through such a limit.
The numbers that admit the sharpest statements---mass and coupling ladders, curvature ratios, shell thermodynamics, and rational witnesses in \texttt{HQIV\_LEAN}---are objects of \emph{discrete mathematics} on $\mathbb{N}$ (combinatorial growth, cumulative simplex ledgers) together with the ordered field $\mathbb{R}$ as used in the proof assistant for analysis, bounds, and phenomenological display.
Smooth-manifold or continuum field notation, when it appears, is a \emph{translation layer} for comparison with standard physics literature; it does not supersede the discrete definitions.

\paragraph{A continuum is not required, and in places would destroy these results.}
The discrete patch layer is \emph{not} a numerical approximation to a more fundamental continuum theory awaiting future construction; it is the ontology.  Where the literature treats ``the continuum limit'' as the goal, HQIV treats it as a \emph{calculation approximation} for comparison with standard formulas.  The continuum form is an optional convenience for comparison; the proved discrete statement is what carries the physics.  In several load-bearing places---the curvature imprint $\delta_E(m)$ at shell $m$, the harmonic ladder masses $m_\tau\to m_\mu\to m_e$, the lock-in vev $v=\sqrt{\eta_{\rm paper}\,\Omega_k(m_{\rm lockin};m_{\rm lockin})}$, the proton anchor at $m=\mathrm{referenceM}=4$, the baryon-asymmetry $\eta\propto 6^7\sqrt{3}/(m+1)$, and the rational coefficients $\alpha=3/5$, $\gamma=2/5$, $\alpha_{\rm GUT}=1/42$---taking a literal $m\to\infty$ or sub-Planck continuum limit would \emph{erase} the discrete index that is doing the predictive work.  Continuum forms of these objects are therefore not preferred refinements; they are degenerate, information-losing readouts of the same patch data.

\paragraph{An exception: $\mathfrak{so}(8)$ as a de-facto continuum algebra from a discrete construction.}
Principled exception to the discrete-only stance: the closed Lie algebra $\mathfrak{so}(8) \supseteq G_2\cup\{\Delta\}$ generated from the Fano octonion left-multiplication table and the phase-lift generator $\Delta\in(\mathrm{e}_1,\mathrm{e}_7)$.  Although $\mathfrak{so}(8)$ is a continuous (28-dimensional real) Lie algebra, it is built here entirely from \emph{discrete} ingredients (the seven imaginary octonion units, a finite Fano table, and one $\pi/2$ rotation), and its closure to 28 dimensions is a finite linear-algebra certificate (see the \texttt{Hqiv.SO8Closure} / \texttt{Hqiv.SO8ClosureInterface} stack).  The Lie-group structure is therefore a \emph{de-facto continuum algebra obtained from a discrete construction}: continuous in parameter space, but with no continuum spacetime, no continuum measure, and no UV limit attached.  When this manuscript or its companions write $\exp(t\,\Delta)\in\mathrm{SO}(8)$ or treat $\mathfrak{so}(8)$ rotations as smooth, those are statements internal to a finite-dimensional Lie group whose generators are certified by discrete data; they are \emph{not} a continuum-spacetime commitment and do not require a continuum-QFT scaffolding.


\paragraph{Patch QFT and discrete numerics (reader contract).}
HQIV (Horizon-Quantized Informational Vacuum) is formulated as a \emph{patch theory}: quantum fields, actions, and physical readouts are attached to discrete patches---shell indices $m\in\mathbb{N}$, finite spacetime index types (e.g.\ $\mathrm{Fin}\,4$), and horizon-limited charts---rather than to a hypothetical fundamental smooth spacetime obtained as a mesh-refinement or ultraviolet limit.
Accordingly, when this text refers to ``QFT,'' it means \emph{patch QFT}: patching rules, finite measurement ledgers, and variational identities on the discrete spine; it is not the claim that the theory is \emph{defined} by recovering ordinary continuum quantum field theory through such a limit.
The numbers that admit the sharpest statements---mass and coupling ladders, curvature ratios, shell thermodynamics, and rational witnesses in \texttt{HQIV\_LEAN}---are objects of \emph{discrete mathematics} on $\mathbb{N}$ (combinatorial growth, cumulative simplex ledgers) together with the ordered field $\mathbb{R}$ as used in the proof assistant for analysis, bounds, and phenomenological display.
Smooth-manifold or continuum field notation, when it appears, is a \emph{translation layer} for comparison with standard physics literature; it does not supersede the discrete definitions.

\paragraph{A continuum is not required, and in places would destroy these results.}
The discrete patch layer is \emph{not} a numerical approximation to a more fundamental continuum theory awaiting future construction; it is the ontology.  Where the literature treats ``the continuum limit'' as the goal, HQIV treats it as a \emph{calculation approximation} for comparison with standard formulas.  The continuum form is an optional convenience for comparison; the proved discrete statement is what carries the physics.  In several load-bearing places---the curvature imprint $\delta_E(m)$ at shell $m$, the harmonic ladder masses $m_\tau\to m_\mu\to m_e$, the lock-in vev $v=\sqrt{\eta_{\rm paper}\,\Omega_k(m_{\rm lockin};m_{\rm lockin})}$, the proton anchor at $m=\mathrm{referenceM}=4$, the baryon-asymmetry $\eta\propto 6^7\sqrt{3}/(m+1)$, and the rational coefficients $\alpha=3/5$, $\gamma=2/5$, $\alpha_{\rm GUT}=1/42$---taking a literal $m\to\infty$ or sub-Planck continuum limit would \emph{erase} the discrete index that is doing the predictive work.  Continuum forms of these objects are therefore not preferred refinements; they are degenerate, information-losing readouts of the same patch data.

\paragraph{An exception: $\mathfrak{so}(8)$ as a de-facto continuum algebra from a discrete construction.}
Principled exception to the discrete-only stance: the closed Lie algebra $\mathfrak{so}(8) \supseteq G_2\cup\{\Delta\}$ generated from the Fano octonion left-multiplication table and the phase-lift generator $\Delta\in(\mathrm{e}_1,\mathrm{e}_7)$.  Although $\mathfrak{so}(8)$ is a continuous (28-dimensional real) Lie algebra, it is built here entirely from \emph{discrete} ingredients (the seven imaginary octonion units, a finite Fano table, and one $\pi/2$ rotation), and its closure to 28 dimensions is a finite linear-algebra certificate (see the \texttt{Hqiv.SO8Closure} / \texttt{Hqiv.SO8ClosureInterface} stack).  The Lie-group structure is therefore a \emph{de-facto continuum algebra obtained from a discrete construction}: continuous in parameter space, but with no continuum spacetime, no continuum measure, and no UV limit attached.  When this manuscript or its companions write $\exp(t\,\Delta)\in\mathrm{SO}(8)$ or treat $\mathfrak{so}(8)$ rotations as smooth, those are statements internal to a finite-dimensional Lie group whose generators are certified by discrete data; they are \emph{not} a continuum-spacetime commitment and do not require a continuum-QFT scaffolding.

% From papers/paper/*.tex: % Shared reader contract: patch QFT + discrete numerics.
% From papers/*.tex:     % Shared reader contract: patch QFT + discrete numerics.
% From papers/*.tex:     \input{include/patch_theory_messaging.tex}
% From papers/paper/*.tex: \input{../include/patch_theory_messaging.tex}

\paragraph{Patch QFT and discrete numerics (reader contract).}
HQIV (Horizon-Quantized Informational Vacuum) is formulated as a \emph{patch theory}: quantum fields, actions, and physical readouts are attached to discrete patches---shell indices $m\in\mathbb{N}$, finite spacetime index types (e.g.\ $\mathrm{Fin}\,4$), and horizon-limited charts---rather than to a hypothetical fundamental smooth spacetime obtained as a mesh-refinement or ultraviolet limit.
Accordingly, when this text refers to ``QFT,'' it means \emph{patch QFT}: patching rules, finite measurement ledgers, and variational identities on the discrete spine; it is not the claim that the theory is \emph{defined} by recovering ordinary continuum quantum field theory through such a limit.
The numbers that admit the sharpest statements---mass and coupling ladders, curvature ratios, shell thermodynamics, and rational witnesses in \texttt{HQIV\_LEAN}---are objects of \emph{discrete mathematics} on $\mathbb{N}$ (combinatorial growth, cumulative simplex ledgers) together with the ordered field $\mathbb{R}$ as used in the proof assistant for analysis, bounds, and phenomenological display.
Smooth-manifold or continuum field notation, when it appears, is a \emph{translation layer} for comparison with standard physics literature; it does not supersede the discrete definitions.

\paragraph{A continuum is not required, and in places would destroy these results.}
The discrete patch layer is \emph{not} a numerical approximation to a more fundamental continuum theory awaiting future construction; it is the ontology.  Where the literature treats ``the continuum limit'' as the goal, HQIV treats it as a \emph{calculation approximation} for comparison with standard formulas.  The continuum form is an optional convenience for comparison; the proved discrete statement is what carries the physics.  In several load-bearing places---the curvature imprint $\delta_E(m)$ at shell $m$, the harmonic ladder masses $m_\tau\to m_\mu\to m_e$, the lock-in vev $v=\sqrt{\eta_{\rm paper}\,\Omega_k(m_{\rm lockin};m_{\rm lockin})}$, the proton anchor at $m=\mathrm{referenceM}=4$, the baryon-asymmetry $\eta\propto 6^7\sqrt{3}/(m+1)$, and the rational coefficients $\alpha=3/5$, $\gamma=2/5$, $\alpha_{\rm GUT}=1/42$---taking a literal $m\to\infty$ or sub-Planck continuum limit would \emph{erase} the discrete index that is doing the predictive work.  Continuum forms of these objects are therefore not preferred refinements; they are degenerate, information-losing readouts of the same patch data.

\paragraph{An exception: $\mathfrak{so}(8)$ as a de-facto continuum algebra from a discrete construction.}
Principled exception to the discrete-only stance: the closed Lie algebra $\mathfrak{so}(8) \supseteq G_2\cup\{\Delta\}$ generated from the Fano octonion left-multiplication table and the phase-lift generator $\Delta\in(\mathrm{e}_1,\mathrm{e}_7)$.  Although $\mathfrak{so}(8)$ is a continuous (28-dimensional real) Lie algebra, it is built here entirely from \emph{discrete} ingredients (the seven imaginary octonion units, a finite Fano table, and one $\pi/2$ rotation), and its closure to 28 dimensions is a finite linear-algebra certificate (see the \texttt{Hqiv.SO8Closure} / \texttt{Hqiv.SO8ClosureInterface} stack).  The Lie-group structure is therefore a \emph{de-facto continuum algebra obtained from a discrete construction}: continuous in parameter space, but with no continuum spacetime, no continuum measure, and no UV limit attached.  When this manuscript or its companions write $\exp(t\,\Delta)\in\mathrm{SO}(8)$ or treat $\mathfrak{so}(8)$ rotations as smooth, those are statements internal to a finite-dimensional Lie group whose generators are certified by discrete data; they are \emph{not} a continuum-spacetime commitment and do not require a continuum-QFT scaffolding.

% From papers/paper/*.tex: % Shared reader contract: patch QFT + discrete numerics.
% From papers/*.tex:     \input{include/patch_theory_messaging.tex}
% From papers/paper/*.tex: \input{../include/patch_theory_messaging.tex}

\paragraph{Patch QFT and discrete numerics (reader contract).}
HQIV (Horizon-Quantized Informational Vacuum) is formulated as a \emph{patch theory}: quantum fields, actions, and physical readouts are attached to discrete patches---shell indices $m\in\mathbb{N}$, finite spacetime index types (e.g.\ $\mathrm{Fin}\,4$), and horizon-limited charts---rather than to a hypothetical fundamental smooth spacetime obtained as a mesh-refinement or ultraviolet limit.
Accordingly, when this text refers to ``QFT,'' it means \emph{patch QFT}: patching rules, finite measurement ledgers, and variational identities on the discrete spine; it is not the claim that the theory is \emph{defined} by recovering ordinary continuum quantum field theory through such a limit.
The numbers that admit the sharpest statements---mass and coupling ladders, curvature ratios, shell thermodynamics, and rational witnesses in \texttt{HQIV\_LEAN}---are objects of \emph{discrete mathematics} on $\mathbb{N}$ (combinatorial growth, cumulative simplex ledgers) together with the ordered field $\mathbb{R}$ as used in the proof assistant for analysis, bounds, and phenomenological display.
Smooth-manifold or continuum field notation, when it appears, is a \emph{translation layer} for comparison with standard physics literature; it does not supersede the discrete definitions.

\paragraph{A continuum is not required, and in places would destroy these results.}
The discrete patch layer is \emph{not} a numerical approximation to a more fundamental continuum theory awaiting future construction; it is the ontology.  Where the literature treats ``the continuum limit'' as the goal, HQIV treats it as a \emph{calculation approximation} for comparison with standard formulas.  The continuum form is an optional convenience for comparison; the proved discrete statement is what carries the physics.  In several load-bearing places---the curvature imprint $\delta_E(m)$ at shell $m$, the harmonic ladder masses $m_\tau\to m_\mu\to m_e$, the lock-in vev $v=\sqrt{\eta_{\rm paper}\,\Omega_k(m_{\rm lockin};m_{\rm lockin})}$, the proton anchor at $m=\mathrm{referenceM}=4$, the baryon-asymmetry $\eta\propto 6^7\sqrt{3}/(m+1)$, and the rational coefficients $\alpha=3/5$, $\gamma=2/5$, $\alpha_{\rm GUT}=1/42$---taking a literal $m\to\infty$ or sub-Planck continuum limit would \emph{erase} the discrete index that is doing the predictive work.  Continuum forms of these objects are therefore not preferred refinements; they are degenerate, information-losing readouts of the same patch data.

\paragraph{An exception: $\mathfrak{so}(8)$ as a de-facto continuum algebra from a discrete construction.}
Principled exception to the discrete-only stance: the closed Lie algebra $\mathfrak{so}(8) \supseteq G_2\cup\{\Delta\}$ generated from the Fano octonion left-multiplication table and the phase-lift generator $\Delta\in(\mathrm{e}_1,\mathrm{e}_7)$.  Although $\mathfrak{so}(8)$ is a continuous (28-dimensional real) Lie algebra, it is built here entirely from \emph{discrete} ingredients (the seven imaginary octonion units, a finite Fano table, and one $\pi/2$ rotation), and its closure to 28 dimensions is a finite linear-algebra certificate (see the \texttt{Hqiv.SO8Closure} / \texttt{Hqiv.SO8ClosureInterface} stack).  The Lie-group structure is therefore a \emph{de-facto continuum algebra obtained from a discrete construction}: continuous in parameter space, but with no continuum spacetime, no continuum measure, and no UV limit attached.  When this manuscript or its companions write $\exp(t\,\Delta)\in\mathrm{SO}(8)$ or treat $\mathfrak{so}(8)$ rotations as smooth, those are statements internal to a finite-dimensional Lie group whose generators are certified by discrete data; they are \emph{not} a continuum-spacetime commitment and do not require a continuum-QFT scaffolding.


\paragraph{Patch QFT and discrete numerics (reader contract).}
HQIV (Horizon-Quantized Informational Vacuum) is formulated as a \emph{patch theory}: quantum fields, actions, and physical readouts are attached to discrete patches---shell indices $m\in\mathbb{N}$, finite spacetime index types (e.g.\ $\mathrm{Fin}\,4$), and horizon-limited charts---rather than to a hypothetical fundamental smooth spacetime obtained as a mesh-refinement or ultraviolet limit.
Accordingly, when this text refers to ``QFT,'' it means \emph{patch QFT}: patching rules, finite measurement ledgers, and variational identities on the discrete spine; it is not the claim that the theory is \emph{defined} by recovering ordinary continuum quantum field theory through such a limit.
The numbers that admit the sharpest statements---mass and coupling ladders, curvature ratios, shell thermodynamics, and rational witnesses in \texttt{HQIV\_LEAN}---are objects of \emph{discrete mathematics} on $\mathbb{N}$ (combinatorial growth, cumulative simplex ledgers) together with the ordered field $\mathbb{R}$ as used in the proof assistant for analysis, bounds, and phenomenological display.
Smooth-manifold or continuum field notation, when it appears, is a \emph{translation layer} for comparison with standard physics literature; it does not supersede the discrete definitions.

\paragraph{A continuum is not required, and in places would destroy these results.}
The discrete patch layer is \emph{not} a numerical approximation to a more fundamental continuum theory awaiting future construction; it is the ontology.  Where the literature treats ``the continuum limit'' as the goal, HQIV treats it as a \emph{calculation approximation} for comparison with standard formulas.  The continuum form is an optional convenience for comparison; the proved discrete statement is what carries the physics.  In several load-bearing places---the curvature imprint $\delta_E(m)$ at shell $m$, the harmonic ladder masses $m_\tau\to m_\mu\to m_e$, the lock-in vev $v=\sqrt{\eta_{\rm paper}\,\Omega_k(m_{\rm lockin};m_{\rm lockin})}$, the proton anchor at $m=\mathrm{referenceM}=4$, the baryon-asymmetry $\eta\propto 6^7\sqrt{3}/(m+1)$, and the rational coefficients $\alpha=3/5$, $\gamma=2/5$, $\alpha_{\rm GUT}=1/42$---taking a literal $m\to\infty$ or sub-Planck continuum limit would \emph{erase} the discrete index that is doing the predictive work.  Continuum forms of these objects are therefore not preferred refinements; they are degenerate, information-losing readouts of the same patch data.

\paragraph{An exception: $\mathfrak{so}(8)$ as a de-facto continuum algebra from a discrete construction.}
Principled exception to the discrete-only stance: the closed Lie algebra $\mathfrak{so}(8) \supseteq G_2\cup\{\Delta\}$ generated from the Fano octonion left-multiplication table and the phase-lift generator $\Delta\in(\mathrm{e}_1,\mathrm{e}_7)$.  Although $\mathfrak{so}(8)$ is a continuous (28-dimensional real) Lie algebra, it is built here entirely from \emph{discrete} ingredients (the seven imaginary octonion units, a finite Fano table, and one $\pi/2$ rotation), and its closure to 28 dimensions is a finite linear-algebra certificate (see the \texttt{Hqiv.SO8Closure} / \texttt{Hqiv.SO8ClosureInterface} stack).  The Lie-group structure is therefore a \emph{de-facto continuum algebra obtained from a discrete construction}: continuous in parameter space, but with no continuum spacetime, no continuum measure, and no UV limit attached.  When this manuscript or its companions write $\exp(t\,\Delta)\in\mathrm{SO}(8)$ or treat $\mathfrak{so}(8)$ rotations as smooth, those are statements internal to a finite-dimensional Lie group whose generators are certified by discrete data; they are \emph{not} a continuum-spacetime commitment and do not require a continuum-QFT scaffolding.


\paragraph{Patch QFT and discrete numerics (reader contract).}
HQIV (Horizon-Quantized Informational Vacuum) is formulated as a \emph{patch theory}: quantum fields, actions, and physical readouts are attached to discrete patches---shell indices $m\in\mathbb{N}$, finite spacetime index types (e.g.\ $\mathrm{Fin}\,4$), and horizon-limited charts---rather than to a hypothetical fundamental smooth spacetime obtained as a mesh-refinement or ultraviolet limit.
Accordingly, when this text refers to ``QFT,'' it means \emph{patch QFT}: patching rules, finite measurement ledgers, and variational identities on the discrete spine; it is not the claim that the theory is \emph{defined} by recovering ordinary continuum quantum field theory through such a limit.
The numbers that admit the sharpest statements---mass and coupling ladders, curvature ratios, shell thermodynamics, and rational witnesses in \texttt{HQIV\_LEAN}---are objects of \emph{discrete mathematics} on $\mathbb{N}$ (combinatorial growth, cumulative simplex ledgers) together with the ordered field $\mathbb{R}$ as used in the proof assistant for analysis, bounds, and phenomenological display.
Smooth-manifold or continuum field notation, when it appears, is a \emph{translation layer} for comparison with standard physics literature; it does not supersede the discrete definitions.

\paragraph{A continuum is not required, and in places would destroy these results.}
The discrete patch layer is \emph{not} a numerical approximation to a more fundamental continuum theory awaiting future construction; it is the ontology.  Where the literature treats ``the continuum limit'' as the goal, HQIV treats it as a \emph{calculation approximation} for comparison with standard formulas.  The continuum form is an optional convenience for comparison; the proved discrete statement is what carries the physics.  In several load-bearing places---the curvature imprint $\delta_E(m)$ at shell $m$, the harmonic ladder masses $m_\tau\to m_\mu\to m_e$, the lock-in vev $v=\sqrt{\eta_{\rm paper}\,\Omega_k(m_{\rm lockin};m_{\rm lockin})}$, the proton anchor at $m=\mathrm{referenceM}=4$, the baryon-asymmetry $\eta\propto 6^7\sqrt{3}/(m+1)$, and the rational coefficients $\alpha=3/5$, $\gamma=2/5$, $\alpha_{\rm GUT}=1/42$---taking a literal $m\to\infty$ or sub-Planck continuum limit would \emph{erase} the discrete index that is doing the predictive work.  Continuum forms of these objects are therefore not preferred refinements; they are degenerate, information-losing readouts of the same patch data.

\paragraph{An exception: $\mathfrak{so}(8)$ as a de-facto continuum algebra from a discrete construction.}
Principled exception to the discrete-only stance: the closed Lie algebra $\mathfrak{so}(8) \supseteq G_2\cup\{\Delta\}$ generated from the Fano octonion left-multiplication table and the phase-lift generator $\Delta\in(\mathrm{e}_1,\mathrm{e}_7)$.  Although $\mathfrak{so}(8)$ is a continuous (28-dimensional real) Lie algebra, it is built here entirely from \emph{discrete} ingredients (the seven imaginary octonion units, a finite Fano table, and one $\pi/2$ rotation), and its closure to 28 dimensions is a finite linear-algebra certificate (see the \texttt{Hqiv.SO8Closure} / \texttt{Hqiv.SO8ClosureInterface} stack).  The Lie-group structure is therefore a \emph{de-facto continuum algebra obtained from a discrete construction}: continuous in parameter space, but with no continuum spacetime, no continuum measure, and no UV limit attached.  When this manuscript or its companions write $\exp(t\,\Delta)\in\mathrm{SO}(8)$ or treat $\mathfrak{so}(8)$ rotations as smooth, those are statements internal to a finite-dimensional Lie group whose generators are certified by discrete data; they are \emph{not} a continuum-spacetime commitment and do not require a continuum-QFT scaffolding.


\section{Scope and claim discipline}\label{sec:scope}

This paper presents a derivation of Kirchhoff's law of thermal
emission inside HQIV.  We do \emph{not} reprove the classical
quantum statistical mechanics; we replace its mode-density input
with HQIV's finite stars-and-bars count and show that the
resulting blackbody machinery (i) is finite at every step,
(ii) admits explicit UV and IR cutoffs by construction, and
(iii) supports one downstream quantitative prediction
($\beta$ at present epoch) that survives the data, under HQIV's
own scale-witness discipline.

\begin{enumerate}[leftmargin=2em]
\item \textbf{Two axioms.}  Discrete null lattice with Planck cell
  size; informational monogamy fixing $\alpha=3/5$ via the
  hockey-stick identity.  These two are the only HQIV
  combinatorial inputs.  The closed Lie algebra spanned by
  $G_2\cup\{\Delta\}\Rightarrow\mathfrak{so}(8)$ used downstream is
  imported as a black-box certificate from the closure paper
  \cite{hqiv-so8-paper}; the rational $(\alpha,\gamma)=(3/5,2/5)$
  pair used here is forced by the same monogamy ledger that the
  3D causal-growth paper \cite{hqiv-3d-growth-paper} derives
  combinatorially.
\item \textbf{Finite by construction.}  All spectral sums are
  $\sum_{m\in[m_{\rm UV},m_{\rm IR}]}$ over finite shell intervals
  in $\mathbb{N}$.  No limits, no analytic continuation, no
  cutoff-removal procedure.  In particular, no programme of
  taking a continuum mesh-refinement limit is invoked or
  required: per the patch-theory reader contract above, several
  load-bearing results would be erased rather than refined by
  such a limit (the rational $\alpha$, the integer cycle index,
  the discrete ladder that the predictive content runs on).
\item \textbf{Single active scale witness.}  HQIV's standing
  discipline (\texttt{AGENTS/ASSUMPTIONS.md} \S7g of the
  HQIV-LEAN repository) admits exactly one dimensionful witness
  per solve.  Throughout this paper we use the default
  \hyperref[lean:proton-lockin]{proton lock-in scale-witness}:
  the proton mass at the lock-in shell $m=4$ uses the theory's
  internal constituent-plus-binding readout, with the observed
  value $M_p=938.272\,\mathrm{MeV}$ serving as the active scale
  witness that pins the mass/unit chart.  The shell index $m=4$
  is the QCD-onset identification (not a free parameter), and
  the constituent-plus-binding closure at that shell is a
  non-trivial HQIV calculation, but the numerical scale of
  $938.272\,\mathrm{MeV}$ is the \emph{anchor}, not an independent
  two-axiom output.  See \S\ref{sec:proton} for the honest
  accounting.
\item \textbf{Wall-clock age from the HQIV-modified CLASS
  implementation.}  The age $t_{\rm wall}=51.2\,\mathrm{Gyr}$
  used in the $\beta$ prediction is an output of the
  HQIV-modified CLASS implementation (a discrete null-lattice
  background with the ADM lapse $N(t)=1+\Phi+\varphi t$ and
  $\varphi$-dependent evolution; the
  \hyperref[lean:wall-clock]{HQVM wall-clock compression}), not
  a free parameter.  Appendix~\ref{app:lean-index} carries the
  Lean snapshot pointer.
\item \textbf{Predictions vs anchor.}  Under the proton
  lock-in scale-witness, observables such as CODATA $1/\alpha$,
  cosmic birefringence $\beta$, and CMB horizon scales are
  \emph{predictions or comparison layers}, not simultaneous
  anchors.  This paper exhibits $\beta$ as one such prediction.
\item \textbf{$\beta$ is non-adjustable under the chosen
  witness.}  Given the anchor, $\beta_{\rm HQIV}$ uses only the
  HQIV stars-and-bars count, the HQIV wall-clock age from the
  HQIV-modified CLASS implementation, and FIRAS $T_{\rm CMB}$
  \cite{fixsen2009}.  No birefringence-side anchor and no
  Friedmann equation enter the calculation.  The internal
  discriminator wall-clock vs apparent age is a single-shot test.
\end{enumerate}

\section{Background: classical Kirchhoff and the normalisation
problem}\label{sec:background}

Kirchhoff's law of thermal emission \cite{kirchhoff1860} states
that, at thermodynamic equilibrium, the spectral radiance emitted
per mode equals the spectral radiance absorbed per mode---the
universal blackbody spectrum.  Two ingredients enter:
\begin{enumerate}[leftmargin=2em]
\item Detailed balance per mode (microscopic reversibility);
\item A mode-density per frequency interval $dN/d\nu$.
\end{enumerate}
In the standard derivation, (i) is statistical-mechanical and
robust; (ii) is the problematic input.  Rayleigh--Jeans gives
$dN/d\nu\propto\nu^{2}$ in a 3D cavity, unbounded above, and so
the classical mean-energy multiplication $k_{\rm B}T\cdot dN/d\nu$
diverges as $\nu\to\infty$---the ultraviolet catastrophe.
Planck's quantum cure \cite{planck1901} provides
$\langle E\rangle=h\nu/(e^{h\nu/k_{\rm B}T}-1)$, which
exponentially suppresses high-$\nu$ contributions and yields a
finite total energy.

This works numerically.  But the cure is asymmetric:
\emph{the high-frequency cutoff is contributed by occupation
suppression, not by mode counting}.  In particular,
nothing about the classical derivation prevents the existence of
modes at arbitrarily high frequency; one simply trusts the Bose
factor to render them harmless.  This is the source of the
normalisation question --- if the mode count itself is unbounded,
any quantum-gravity or trans-Planckian completion has to act on
the mode list, not just on the occupation.

HQIV makes the mode count finite at the start.

\section{HQIV mode-count axiom}\label{sec:axiom}

\begin{definition}[Discrete null-lattice mode count]\label{def:simplex}
Let $\mathcal{N}(m)$ be the count of new Planck-scale lattice
cells at temperature-ladder depth $m\in\mathbb{N}$.  HQIV's
discrete null-lattice axiom gives
\begin{equation}\label{eq:NM}
  \mathcal{N}(m)=(m+2)(m+1).
\end{equation}
This is the integer-solution count for $x+y+z=m$, $x,y,z\ge 0$,
twice the binomial $\binom{m+2}{2}$.
\end{definition}

\begin{proposition}[Hockey-stick cumulative count]\label{prop:cum}
\begin{equation}\label{eq:cum}
  \sum_{m'=0}^{n}\mathcal{N}(m')
   =\tfrac{1}{3}(n+1)(n+2)(n+3).
\end{equation}
\end{proposition}
\begin{proof}
By induction on $n$, using \eqref{eq:NM}; the
\hyperref[lean:hockey-stick]{hockey-stick cumulative-count
identity} states the equivalent integer form
$3\sum_{m'=0}^{n}\mathcal{N}(m')=(n+1)(n+2)(n+3)$.
\end{proof}

\begin{corollary}[Both cutoffs finite]\label{cor:finite}
For any $0\le m_{\rm UV}\le m_{\rm IR}<\infty$,
\[
  \sum_{m\in[m_{\rm UV},m_{\rm IR}]}\mathcal{N}(m)<\infty.
\]
\end{corollary}

This is the structural fact that lets the entire spectral
machinery be a finite sum, regardless of the temperature and
without any regularisation procedure.

\section{HQIV temperature and frequency ladders}\label{sec:ladder}

The HQIV harmonic ladder fixes the shell-indexed frequency at:
\begin{equation}\label{eq:freq}
  \omega_m=\frac{T_{\rm Pl}}{m+1},\qquad m\in\mathbb{N},
\end{equation}
with $\omega_{0}=T_{\rm Pl}$ at the Planck pole (highest
frequency, no UV catastrophe) and $\omega_m\to 0$ as
$m\to\infty$ (IR limit, only $\mathcal{N}(m)\sim m^{2}$ modes
each).  Equivalently, the temperature ladder is
$T(m)=T_{\rm Pl}/(m+1)$, exposed as the
\hyperref[lean:auxiliary-T]{auxiliary-field temperature ladder}.

\section{Kirchhoff's law on a finite shell list}\label{sec:kirchhoff}

We now state Kirchhoff's law in HQIV form.  Two ingredients,
both finite by construction.

\begin{definition}[Per-shell Bose occupation]\label{def:bose}
At temperature $T>0$ and shell $m\in\mathbb{N}$,
\begin{equation}\label{eq:nbose}
  n_{B}(\omega_m,T)
   =\frac{1}{\exp(\omega_m/T)-1}.
\end{equation}
This is the standard Bose factor on \emph{one} HQIV mode---the
\hyperref[lean:bose]{per-shell Bose factor}. It is positive and
finite at every $m\in\mathbb{N}$ and $T>0$ (with the divergence at
$\omega_m=T\to 0$ outside the shell-discrete sector).
\end{definition}

\begin{definition}[Per-shell spectral energy]\label{def:spec}
\begin{equation}\label{eq:shellE}
  E_m(T)=\mathcal{N}(m)\cdot\omega_m\cdot n_{B}(\omega_m,T).
\end{equation}
Formalised as the
\hyperref[lean:shell-spectral]{per-shell spectral-energy density}.
Every factor is positive at every $m,T$; no regulariser is
applied.
\end{definition}

\begin{definition}[Truncated blackbody energy density]\label{def:bb}
For any $m_{\rm UV}\le m_{\rm IR}$ in $\mathbb{N}$,
\begin{equation}\label{eq:Edens}
  u(T;m_{\rm UV},m_{\rm IR})
   =\sum_{m=m_{\rm UV}}^{m_{\rm IR}} E_m(T).
\end{equation}
Formalised as the
\hyperref[lean:blackbody-density]{truncated blackbody energy
density}.
\end{definition}

\begin{theorem}[Finite-mode Kirchhoff law]\label{thm:kirchhoff}
For any $T>0$ and any $m_{\rm UV}\le m_{\rm IR}$ in $\mathbb{N}$,
\eqref{eq:Edens} is a finite, positive, monotone sum:
\begin{enumerate}[leftmargin=2em]
\item $0\le u(T;m_{\rm UV},m_{\rm IR})<\infty$;
\item $0<u(T;m_{\rm UV},m_{\rm IR})$ whenever $m_{\rm UV}\le m_{\rm IR}$;
\item monotone in both cutoffs: lowering $m_{\rm UV}$ or raising
  $m_{\rm IR}$ adds non-negative contributions.
\end{enumerate}
The result is true without any limiting procedure or
renormalisation: it holds as a closed-form identity on
$\mathbb{N}^{2}\times\mathbb{R}_{>0}$.
\end{theorem}
\begin{proof}
The
\hyperref[lean:blackbody-density]{truncated blackbody energy
density} is non-negative and finite as a finite-range sum of
non-negative terms; it is strictly positive on every nonempty
shell window and monotone in both cutoff arguments. All three
parts are zero-\texttt{sorry} Lean theorems in the same module
(see Appendix~\ref{app:lean-index}).
\end{proof}

\begin{remark}[The UV catastrophe never appears]\label{rmk:nouv}
Unlike the classical derivation, no high-frequency divergence
exists to suppress.  $\mathcal{N}(m)$ is exactly $(m+2)(m+1)$ at
the Planck pole $m=0$, namely 2 modes (lattice points
$(0,0,1),(0,1,0),(1,0,0)$ scaled to two orientations), and there
is no shell above it: $\omega>T_{\rm Pl}$ does not exist in the
HQIV mode list.  The Bose factor in \eqref{eq:nbose} additionally
provides exponential suppression at fixed $T\ll T_{\rm Pl}$, so
the Planck spectrum shape is preserved everywhere classical
spectroscopy probes; but the structural finiteness is independent
of any thermal cutoff.
\end{remark}

\begin{remark}[The IR is also finite]\label{rmk:noir}
At fixed temperature $T$, the Bose factor diverges as
$\omega_m\to 0$ ($m\to\infty$): $n_B(\omega_m,T)\sim T/\omega_m=
T(m+1)/T_{\rm Pl}$, while $\omega_m\to 0$.  The product
$\omega_m\,n_B(\omega_m,T)\to T$, finite per mode; multiplying by
$\mathcal{N}(m)\sim m^{2}$ gives a per-shell contribution growing
as $m^{2}T$.  If one took $m_{\rm IR}\to\infty$ the sum would
diverge, exactly the classical Rayleigh--Jeans IR behavior---but
HQIV's discrete null lattice places a structural cutoff at the
lock-in reference shell $M$ (the
\hyperref[lean:ir-cutoff]{lock-in IR cutoff}), making \emph{the
IR also finite at the lattice level}, not just by stipulation.
\end{remark}

\begin{definition}[Vacuum zero-point cell]\label{def:zp}
\begin{equation}\label{eq:zp}
  u_{0}(m_{\rm UV},m_{\rm IR})
   =\sum_{m=m_{\rm UV}}^{m_{\rm IR}}
    \frac{\mathcal{N}(m)\,\omega_m}{2}.
\end{equation}
The
\hyperref[lean:vacuum-zp]{vacuum zero-point cell} is non-negative
on every shell interval, by the same finite-sum-of-non-negatives
argument used in Theorem~\ref{thm:kirchhoff}.
\end{definition}

\begin{theorem}[Decomposition]\label{thm:decomp}
\begin{equation}\label{eq:decomp}
  u_{\rm tot}(T;m_{\rm UV},m_{\rm IR})
   =u_{0}(m_{\rm UV},m_{\rm IR})
   +u(T;m_{\rm UV},m_{\rm IR}).
\end{equation}
Formalised as the
\hyperref[lean:blackbody-with-vacuum]{blackbody-plus-vacuum
energy density}, with the additive decomposition lemma in the
same module.  No renormalisation is required to handle the
zero-point cell---it is a finite, structural part of the spectrum
on a finite shell list.
\end{theorem}

\paragraph{Summary.}
Equation~\eqref{eq:decomp} is the HQIV statement of Kirchhoff's
law.  The standard Stefan--Boltzmann scaling $u(T;0,M)\sim T^{4}$
and the Wien displacement bound emerge from \eqref{eq:Edens} via
the
\hyperref[lean:stefan]{horizon Stefan--Boltzmann module} and the
\hyperref[lean:wien]{horizon Wien-displacement module}, and rely
only on Theorem~\ref{thm:kirchhoff}.  No regularisation appears
in any step.

\section{The scale anchor: proton mass at the lock-in shell
$m=4$}\label{sec:proton}

We use this section to be explicit about what the proton mass
\emph{is} and is \emph{not} doing in this paper, since careless
phrasing here can over-claim independence.

\paragraph{Status under the proton lock-in scale-witness.}
Per \texttt{AGENTS/ASSUMPTIONS.md} \S7g of the HQIV-LEAN
repository, the HQIV pipeline carries a single active
dimensionful scale witness per solve.  The default is the
\hyperref[lean:proton-lockin]{proton lock-in scale-witness}: the
proton mass at the lock-in shell $m=4$ uses the theory's internal
constituent-plus-binding readout (described below), with the
\emph{observed} value $938.272\,\mathrm{MeV}$ serving as the
active scale witness that pins the mass/unit chart.  Other
quantities---CODATA $1/\alpha$, CMB horizons, PDG centrals
\cite{pdg2024}---are then predictions or comparison layers
against this anchor.  The
\hyperref[lean:scale-witness]{scale-witness enum} carries the
documented rule: ``never wire CODATA $\alpha$ \emph{and} proton
mass \emph{and} a third cosmology anchor into the same solve
without documenting which one is active and which are
predictions.''

\paragraph{Shell index $m=4$: chosen, not derived.}  The
reference shell $m=4$ is the QCD-onset identification: it is the
HQIV temperature-ladder index where colour confinement and the
lock-in chart begin.  The choice is not a free fit parameter---
the hadronic mass content of the universe requires \emph{some}
QCD-onset shell, and $m=4$ is the smallest integer consistent
with the small number of post-QCD steps needed to reach the
nucleon harmonic.  But it is also not a two-axiom output: the
two axioms by themselves do not single out $m=4$.

\paragraph{What \emph{is} non-trivial at $m=4$.}  Given the shell
choice, the readout
\begin{equation}\label{eq:proton}
  M_{p}^{\rm HQIV}
   = M^{\rm const}_{p}\,-\,E^{\rm shared}_{\rm bind}
   = 938.272\,\mathrm{MeV}
\end{equation}
packages the nucleon as constituent-energy minus shared binding,
formalised as the
\hyperref[lean:derived-proton]{derived proton mass} (with the
\hyperref[lean:proton-shared]{shared-harmonics closure} establishing
\eqref{eq:proton}).  The constituent harmonics $M^{\rm const}_{p}$
are dictated by the mode count \eqref{eq:NM} at $m=4$; the
shared-binding energy $E^{\rm shared}_{\rm bind}$ has its own
composite-trace construction.  These combinatorial structures are
genuine HQIV content, and the closure to a single number
$938.272\,\mathrm{MeV}$ is a non-trivial constraint on the
lock-in shell-set.  The
\hyperref[lean:proton-anchor]{proton-anchor mass constant} records
this as the active scale witness.

\paragraph{Why this is not a ``test'' in the same sense as $\beta$.}
At the level of the broader theory, the numerical value of
$M_p^{\rm HQIV}$ is the chart-pinning anchor: it defines the unit
of mass/energy used to interpret every other dimensionful
quantity.  Calling it a ``test'' would conflate two distinct
roles: (i) the narrow constituent-plus-binding readout at the
lock-in shell, which is a genuine HQIV calculation, vs (ii) the
position of the lock-in shell on the integer ladder, which is a
modelling choice in the broader theory.  The honest framing is:
$M_p$ is the \emph{anchor}, established once;
$\beta$ is the \emph{prediction}, evaluated against the anchor
and against the data.

\paragraph{Cross-check status.}  Within HQIV's own discipline,
candidate scale-witness alternatives are recorded
(CODATA $\alpha$ witness, present-CMB witness) and exist in the
\hyperref[lean:scale-witness]{scale-witness enum}.  Switching to a
different witness is a deliberate operation and would re-pin every
dimensionful HQIV output.  For this paper all numerics are
reported under the proton lock-in scale-witness.

\section{Quantitative test: cosmic birefringence $\beta=0.379^\circ$}\label{sec:beta}

This is the section where HQIV makes a single-shot quantitative
prediction under the chosen scale witness.  The proton-mass
anchor of \S\ref{sec:proton} has already pinned the unit chart.
The birefringence calculation then combines the HQIV wall-clock
age $t_{\rm wall}=51.2\,\mathrm{Gyr}$---an output of the
HQIV-modified CLASS implementation (a discrete null-lattice
background with the ADM lapse $N(t)=1+\Phi+\varphi t$ and
$\varphi$-dependent evolution; the
\hyperref[lean:wall-clock]{HQVM wall-clock compression})---with
the same stars-and-bars mode count $\mathcal{N}(m)=(m+2)(m+1)$ of
\eqref{eq:NM} under the propagation-shell identification.  Only
one externally measured dimensionful input enters at this stage:
FIRAS $T_{\rm CMB}=2.7255\,\mathrm{K}$ \cite{fixsen2009}.  No
birefringence-side anchor enters---in particular, no Planck PR4
calibration \cite{planckPR42020}---and no Friedmann equation is
imposed.

The Planck DR4 EB analysis of Eskilt \& Komatsu \cite{eskilt2022}
reports an isotropic cosmic birefringence
\begin{equation}\label{eq:pr4}
  \beta_{\rm obs}=0.342^\circ\pm 0.094^\circ
  \quad (\text{Planck PR4 NPIPE, EB cross-spectrum}),
\end{equation}
a $\sim3.6\sigma$ detection, frequency-independent across the
30--353~GHz Planck sub-bands, isotropic within errors.

\paragraph{Two shell coordinates.}  HQIV carries two genuinely
distinct shell coordinates (\S\ref{sec:axiom},
\S\ref{sec:ladder}, and the
\hyperref[lean:cosm-shell-ladder]{cosmological shell-ladder
module}):
\begin{align}
  m_{T}(T_{\rm obs}) &= T_{\rm Pl}/T_{\rm obs}-1
   \quad\text{(real-valued shell index in the temperature-ladder chart),}\\
  m_{\rm prop} &= \text{count of HQIV propagation cells crossed since the Planck pole.}
\end{align}
A photon's proper time is zero, so it does not accumulate
temperature-ladder positions along its worldline; the cumulative
birefringence formula uses the propagation coordinate.

\paragraph{Propagation-shell coarseness.}  Identifying one HQIV
propagation shell at observation depth $m_T$ with the bundle of
new lattice simplices appearing at that depth,
\begin{equation}\label{eq:bundle}
  \text{Planck cells per propagation shell}=\mathcal{N}(m_T),
\end{equation}
the propagation-shell offset at the observation event is
\begin{equation}\label{eq:mprop}
  m_{\rm prop}=\frac{t_{\rm wall}/t_{\rm Pl}}{\mathcal{N}(m_T)}
              =\frac{t_{\rm wall}/t_{\rm Pl}}{(m_T+2)(m_T+1)}
              \;\approx\;
              \frac{t_{\rm wall}}{t_{\rm Pl}}
              \!\left(\frac{T_{\rm obs}}{T_{\rm Pl}}\right)^{\!2}.
\end{equation}
This is the central use of the same stars-and-bars count
\eqref{eq:NM} that delivered the finite-mode Kirchhoff law in
\S\ref{sec:kirchhoff}; we are not introducing a new HQIV input.

\paragraph{Birefringence formula.}  HQIV's
\hyperref[lean:cumbire]{cumulative birefringence observable},
proven zero-\texttt{sorry}, is
\begin{equation}\label{eq:cumbire}
  \beta(m_{\rm emit},m_{\rm obs})
   =\alpha\,\log\!\left(\frac{m_{\rm obs}+1}{m_{\rm emit}+1}\right),
\end{equation}
with $\alpha=3/5$.  Specialising to the near-pole reading
($m_{\rm emit}=0$ at the Planck pole, $m_{\rm obs}=m_{\rm prop}$):
\begin{equation}\label{eq:cumbireNP}
  \beta_{\rm HQIV}=\alpha\,\log(1+m_{\rm prop}).
\end{equation}

\paragraph{Numerics.}  Inputs:
$T_{\rm Pl}=1.41679\times10^{32}\,\mathrm{K}$,
$T_{\rm CMB}=2.7255\,\mathrm{K}$ (FIRAS \cite{fixsen2009}),
$t_{\rm Pl}=5.391\times10^{-44}\,\mathrm{s}$,
$t_{\rm wall}=51.2\,\mathrm{Gyr}$ (the
\hyperref[lean:wall-clock]{wall-clock-age snapshot} from the
HQIV-modified CLASS background),
$\alpha=3/5$ (informational monogamy).  Then
\[
\begin{array}{rcl}
  m_T &=& 5.198\times10^{31},\\
  \mathcal{N}(m_T) &=& 2.702\times10^{63},\\
  t_{\rm wall}/t_{\rm Pl} &=& 2.997\times10^{61},\\
  m_{\rm prop} &=& 1.109\times10^{-2}.
\end{array}
\]
\begin{equation}\label{eq:final}
  \beta_{\rm HQIV}
   =\tfrac{3}{5}\log(1.01109)
   =6.619\times10^{-3}\,\mathrm{rad}
   =\boxed{0.3792^\circ},
\end{equation}
deviation $-0.40\sigma$ from \eqref{eq:pr4}. The exact arithmetic
chain reproducing every line of this paragraph ships in the
supplementary script bundle (Appendix~\ref{app:scripts}, file
\texttt{birefringence\_calculation.py}).

\paragraph{Falsifier: apparent age fails.}  Substituting the
\hyperref[lean:apparent-age]{HQIV apparent age} of
$13.8\,\mathrm{Gyr}$ for $t_{\rm wall}$:
$m_{\rm prop}^{\rm app}\approx 3.0\times10^{-3}$,
$\beta^{\rm app}\approx 0.103^\circ$, a $-2.55\sigma$ deviation.
The formula \emph{selects} HQIV's wall-clock age over the apparent
age---this is a non-adjustable test of the ADM-lapse subsystem
against an independent observation.

\paragraph{Observer-position character (PR4 frequency
independence).}  By \eqref{eq:cumbireNP} and \eqref{eq:mprop},
$\beta_{\rm HQIV}$ depends on $(t_{\rm wall},T_{\rm obs})$ only,
not on individual photon frequencies.  All photons reaching the
same observer get the same imprint, formalised as the
\hyperref[lean:bire-freq-indep]{frequency-independence statement
for the predicted birefringence}, and exactly the PR4 finding
\cite{eskilt2022} across 30--353~GHz. A per-photon $\omega^{2}$
mechanism would predict $\sim 140\times$ amplification across that
range and is ruled out; HQIV does not predict that. The blackbody
Planck spectrum smears in photon energy; the $\beta$ value does
not.

\paragraph{Isotropy.}  $m_{\rm prop}$ is a scalar.
$\beta_{\rm HQIV}$ has no preferred direction, matching the PR4
isotropy null result---formalised as the
\hyperref[lean:bire-isotropic]{isotropy lemma} for the predicted
birefringence (with the
\hyperref[lean:bire-isotropic-near-pole]{near-pole specialisation}
used in \eqref{eq:cumbireNP}).

\section{Friedmann scaling recovered, not imported}\label{sec:friedmann}

In a radiation-dominated FLRW cosmology $H\propto T^{2}/M_{\rm Pl}$
and so the Hubble time is $t_{H}=1/H\propto M_{\rm Pl}/T^{2}$.
This is famously the quadratic growth law.

In HQIV, identifying one propagation shell with the bundle of new
lattice simplices at the observation depth (the same identification
that delivered \eqref{eq:mprop}) gives the propagation-shell
duration $\sim t_{\rm Pl}\,\mathcal{N}(m_T)\sim t_{\rm Pl}(T_{\rm Pl}/T)^{2}
=M_{\rm Pl}/T^{2}$ in Planck units.

\begin{theorem}[Friedmann scaling recovered]\label{thm:friedmann}
Under the propagation-shell identification of \S\ref{sec:beta},
\[
  t_{H}^{\rm HQIV}(T)\sim t_{\rm Pl}\,\mathcal{N}(m_T)
                    \sim M_{\rm Pl}/T^{2}.
\]
\end{theorem}
\noindent\emph{The same stars-and-bars count \eqref{eq:NM} that
delivers the finite-mode Kirchhoff law also delivers Friedmann's
$H\propto T^{2}$.}  No Einstein equation has been used. The
arithmetic ratio $t_H^{\rm HQIV}(T)/t_H^{\rm rad}(T)
=(m_T+2)/(m_T+1)\to 1+O(1/m_T)$ at present-epoch temperatures is
exhibited by the supplementary script
\texttt{friedmann\_recovery.py} described in
Appendix~\ref{app:scripts}.

\section{Discussion: where normalisation went}\label{sec:discussion}

In the standard derivation of the Planck spectrum, a UV
regulariser is implicit (a cavity cutoff, dimensional reg,
or simply trust in the Bose factor); an IR regulariser is also
implicit (the cavity volume).  Renormalisation theory then
constructs observables independent of these regularisers.  This
is conceptually expensive---the cutoffs are tools, not physics.

HQIV trades this for two combinatorial inputs.  The discrete null
lattice provides a structural UV cutoff at $m=0$ (Planck pole);
the lock-in reference shell provides a structural IR cutoff at
$m=M$.  Both cutoffs are physical choices, not regulators: moving
them changes observables in a calculable way, captured by the
monotonicity statements of Theorem~\ref{thm:kirchhoff}.  The
total mode budget at any pair of cutoffs is the closed-form
hockey-stick sum \eqref{eq:cum}.

What is gained:
\begin{itemize}[leftmargin=2em]
\item No UV catastrophe---the high-frequency modes do not exist.
\item No IR divergence---the low-frequency modes are finite in
  count.
\item No renormalisation procedure---the sums are finite as
  written.
\item Under the proton lock-in scale-witness convention, one
  non-trivial quantitative prediction ($\beta=0.379^\circ$ at
  present epoch) survives the data at $-0.40\sigma$ and
  discriminates wall-clock from apparent age internally.
\item Friedmann scaling recovered as a consequence
  (Theorem~\ref{thm:friedmann}), not imported.
\end{itemize}

What is \emph{not} claimed:
\begin{itemize}[leftmargin=2em]
\item HQIV is not claimed to be independent of every scale
  anchor.  The proton mass at $m=4$ uses the theory's internal
  constituent-plus-binding readout, with the observed value
  $938.272\,\mathrm{MeV}$ serving as the active scale witness.
  Substituting an alternative witness (CODATA $\alpha$,
  present-CMB) would re-pin every dimensionful HQIV output;
  no such cross-witness consistency demonstration is given here.
\item The shell index $m=4$ is not a two-axiom output.  It is
  the QCD-onset identification, and the constituent-plus-binding
  closure at that shell to $938.272\,\mathrm{MeV}$ is a non-trivial
  HQIV calculation, but the placement of $m=4$ itself is a
  modelling choice.
\item The wall-clock age $51.2\,\mathrm{Gyr}$ is an output of
  the HQIV-modified CLASS implementation (a discrete null-lattice
  background with the ADM lapse and $\varphi$-dependent
  evolution), not a free parameter, but neither is it strictly
  forced by the two axioms alone---the HQIV-modified CLASS
  background carries its own modelling choices on top of the
  lattice.
\item One further piece of the $\beta$ chain is an explicit
  identification, not a strict consequence of the two axioms:
  ``one propagation shell $=$ one bundle of new lattice
  simplices at the observation depth'' (the same construction
  that gives Friedmann's $T^{2}$ scaling).
\item No continuum-limit programme is invoked, required, or
  desired anywhere in this chain. Per the patch-theory reader
  contract, the discrete patch layer is the ontology; the only
  continuum object that appears is the de-facto continuum
  algebra $\mathfrak{so}(8)$ obtained from finite octonionic
  data \cite{hqiv-so8-paper,hqiv-3d-growth-paper}, and even
  that is built from discrete ingredients with no continuum
  spacetime attached.
\end{itemize}

The proper reading of this paper is therefore: HQIV's
combinatorial mode-count axiom delivers Kirchhoff's law on a
finite shell list without any regularisation; under the standing
proton lock-in scale-witness convention, combining the HQIV
wall-clock age from the HQIV-modified CLASS implementation with
the lattice mode count $\mathcal{N}(m)$ under the
propagation-shell identification predicts
$\beta_{\rm HQIV}=0.379^\circ$, surviving the PR4 data
\cite{eskilt2022,planckPR42020} to $-0.40\sigma$.  Stronger
claims of full independence require either a derivation of $m=4$
and $t_{\rm wall}$ from the two axioms or a cross-witness
consistency demonstration; both are open theoretical targets.

\appendix

\section{Lean module index}
\label{app:lean-index}

The Lean development lives in the public slimmed mirror
\url{https://github.com/HQIV/hqiv-lean} \cite{hqiv-lean}.  All
theorems cited in the body are sorry-free at the time of
writing.  Each entry below carries a reader-friendly anchor (used
by hyperlinks in the body), a one-line description, and clickable
links to the Lean module and identifiers.

\begingroup
\sloppy
\setlength{\emergencystretch}{4em}
\begin{description}[leftmargin=0em,style=nextline]
\item[\phantomsection\label{lean:simplex-count}\textbf{Lattice simplex mode count.}]
  The integer-solution count
  $\mathcal{N}(m)=(m+2)(m+1)$ at temperature-ladder depth $m$.
  Lean: \leanid{latticeSimplexCount} and
  \leanid{latticeSimplexCount_eq} in
  \leanfile{Hqiv/Geometry/OctonionicLightCone.lean};
  positivity is recorded as
  \leanid{latticeSimplexCount_pos}.
\item[\phantomsection\label{lean:hockey-stick}\textbf{Hockey-stick cumulative-count identity.}]
  The closed-form sum
  $3\sum_{m'=0}^{n}\mathcal{N}(m')=(n+1)(n+2)(n+3)$.
  Lean: \leanid{cumLatticeSimplexCount_hockey_stick} and
  \leanid{cumLatticeSimplexCount_closed} in
  \leanfile{Hqiv/Geometry/OctonionicLightCone.lean}.
\item[\phantomsection\label{lean:auxiliary-T}\textbf{Auxiliary-field temperature ladder.}]
  $T(m)=T_{\rm Pl}/(m+1)$ on the discrete shell index $m$.
  Lean: \leanid{T} in
  \leanfile{Hqiv/Geometry/AuxiliaryField.lean}, with
  \leanid{shellOmega} (the harmonic ladder of \eqref{eq:freq})
  in
  \leanfile{Hqiv/Physics/HorizonBlackbodySpectrum.lean}.
\item[\phantomsection\label{lean:bose}\textbf{Per-shell Bose factor.}]
  The standard occupation $n_B(\omega_m,T)$ of \eqref{eq:nbose}.
  Lean: \leanid{nBose} (positivity \leanid{nBose_pos}) in
  \leanfile{Hqiv/Physics/HorizonBlackbodySpectrum.lean}.
\item[\phantomsection\label{lean:shell-spectral}\textbf{Per-shell spectral-energy density.}]
  Equation \eqref{eq:shellE},
  $E_m(T)=\mathcal{N}(m)\,\omega_m\,n_B(\omega_m,T)$. Lean:
  \leanid{shellSpectralEnergy} (and
  \leanid{shellModeMultiplicity} for $\mathcal{N}(m)$ in the
  new-modes basis) in
  \leanfile{Hqiv/Physics/HorizonBlackbodySpectrum.lean}.
\item[\phantomsection\label{lean:blackbody-density}\textbf{Truncated blackbody energy density.}]
  Equation \eqref{eq:Edens}, with the three-clause finiteness /
  positivity / monotonicity package. Lean:
  \leanid{blackbodyEnergyDensity},
  \leanid{blackbodyEnergyDensity_pos_of_le},
  \leanid{blackbodyEnergyDensity_mono_UV},
  \leanid{blackbodyEnergyDensity_mono_IR},
  \leanid{planckUVCutoff} (Planck pole) in
  \leanfile{Hqiv/Physics/HorizonBlackbodySpectrum.lean}.
\item[\phantomsection\label{lean:ir-cutoff}\textbf{Lock-in IR cutoff.}]
  The structural infrared cutoff at the lock-in reference shell
  $M$, packaged as
  \leanid{referenceIRCutoff} in
  \leanfile{Hqiv/Physics/HorizonBlackbodySpectrum.lean}.
\item[\phantomsection\label{lean:vacuum-zp}\textbf{Vacuum zero-point cell.}]
  Equation \eqref{eq:zp}; Lean:
  \leanid{vacuumZeroPointEnergy} (non-negativity
  \leanid{vacuumZeroPointEnergy_nonneg}) in
  \leanfile{Hqiv/Physics/HorizonBlackbodySpectrum.lean}.
\item[\phantomsection\label{lean:blackbody-with-vacuum}\textbf{Blackbody-plus-vacuum energy density.}]
  Equation \eqref{eq:decomp}; Lean:
  \leanid{blackbodyEnergyDensityWithVacuum} (and the additive
  decomposition lemma) in
  \leanfile{Hqiv/Physics/HorizonBlackbodySpectrum.lean}.
\item[\phantomsection\label{lean:stefan}\textbf{Horizon Stefan--Boltzmann module.}]
  Recovery of the $T^4$ scaling and Stefan-type ceiling on a
  finite shell list. Lean:
  \leanfile{Hqiv/Physics/HorizonBlackbodyStefan.lean}.
\item[\phantomsection\label{lean:wien}\textbf{Horizon Wien-displacement module.}]
  Wien displacement bound from the same finite-mode spectrum.
  Lean:
  \leanfile{Hqiv/Physics/HorizonBlackbodyWienDisplacement.lean}
  (with the polarized-greybody companion in
  \leanfile{Hqiv/Physics/HorizonBlackbodyGreybody.lean}).
\item[\phantomsection\label{lean:proton-lockin}\textbf{Proton lock-in scale-witness.}]
  The HQIV scale-witness convention pinning the mass/unit chart
  at $m=4$ to $M_p=938.272\,\mathrm{MeV}$. Lean:
  \leanid{ScaleWitness.proton_lockin} (and the
  \hyperref[lean:scale-witness]{scale-witness enum}) in
  \leanfile{Hqiv/Physics/ScaleWitness.lean}.
\item[\phantomsection\label{lean:scale-witness}\textbf{Scale-witness enum.}]
  The Lean enum recording the candidate single-witness anchors
  (\leanid{proton_lockin}, \leanid{codata_alpha},
  \leanid{cmb_now}). Lean: \leanid{ScaleWitness} in
  \leanfile{Hqiv/Physics/ScaleWitness.lean}.
\item[\phantomsection\label{lean:derived-proton}\textbf{Derived proton mass.}]
  The constituent-minus-binding readout
  $M_p^{\rm HQIV}=M_p^{\rm const}-E^{\rm shared}_{\rm bind}$ of
  \eqref{eq:proton}. Lean: \leanid{derivedProtonMass}
  (positivity \leanid{derivedProtonMass_pos}) in
  \leanfile{Hqiv/Physics/DerivedNucleonMass.lean}.
\item[\phantomsection\label{lean:proton-shared}\textbf{Shared-harmonics proton-mass closure.}]
  The closure lemma certifying that the constituent and shared
  pieces sum to $938.272\,\mathrm{MeV}$. Lean:
  \leanid{proton_mass_from_shared_harmonics} in
  \leanfile{Hqiv/Physics/DerivedNucleonMass.lean}.
\item[\phantomsection\label{lean:proton-anchor}\textbf{Proton-anchor mass constant.}]
  The $938.272\,\mathrm{MeV}$ anchor numeral that pins the unit
  chart. Lean: \leanid{protonAnchorMass_MeV} (also
  \leanid{protonMassFromMetaHarmonics_MeV}) in
  \leanfile{Hqiv/Physics/QuarkMetaResonance.lean}.
\item[\phantomsection\label{lean:cosm-shell-ladder}\textbf{Cosmological shell-ladder module.}]
  The propagation-shell coordinate $m_{\rm prop}$ of
  \eqref{eq:mprop}, the simplex-count cast
  \leanid{latticeSimplexCount_as_shell_coarseness}, and the
  internal-coordinate identity
  \leanid{m_prop_HQIV_internal_eq}. Lean:
  \leanfile{Hqiv/Cosmology/CosmologicalShellLadder.lean}.
\item[\phantomsection\label{lean:cumbire}\textbf{Cumulative birefringence observable.}]
  Equation \eqref{eq:cumbire}; the near-pole specialisation
  \eqref{eq:cumbireNP} is
  \leanid{NearPoleObservation.predictedBirefringence}, and the
  numerical pass against PR4 errors is
  \leanid{nearPole_cmb_shell_ladder_pass}. Lean:
  \leanfile{Hqiv/Cosmology/CosmologicalShellLadder.lean} and
  \leanfile{Hqiv/Physics/HorizonBlackbodyLadder.lean}.
\item[\phantomsection\label{lean:bire-isotropic}\textbf{Isotropy of the predicted birefringence.}]
  Lean: \leanid{predictedBirefringence_isotropic} (with the
  \hyperref[lean:bire-isotropic-near-pole]{near-pole
  specialisation}) in
  \leanfile{Hqiv/Cosmology/CosmologicalShellLadder.lean}; the
  no-shell-crossing zero is
  \leanid{predictedBirefringence_zero_of_no_shell_crossing}, and
  the lab-scale upper bound is
  \leanid{nearPole_predictedBirefringence_lab_bound}.
\item[\phantomsection\label{lean:bire-isotropic-near-pole}\textbf{Near-pole isotropy specialisation.}]
  Lean:
  \leanid{nearPole_predictedBirefringence_isotropic} in
  \leanfile{Hqiv/Cosmology/CosmologicalShellLadder.lean}.
\item[\phantomsection\label{lean:bire-freq-indep}\textbf{Frequency-independence of the predicted birefringence.}]
  All photons reaching the same observer carry the same
  birefringence imprint. Lean:
  \leanid{predictedBirefringence_frequency_independent} in
  \leanfile{Hqiv/Cosmology/CosmologicalShellLadder.lean}.
\item[\phantomsection\label{lean:wall-clock}\textbf{HQVM wall-clock compression.}]
  The HQIV-modified CLASS background output recording
  $t_{\rm wall}=51.2\,\mathrm{Gyr}$ for downstream proof reuse.
  Lean: \leanid{age_wall_clock_Gyr_paper} (with
  \leanid{wallClockAgeHomogeneous} and \leanid{HQVM_lapse}) in
  \leanfile{Hqiv/Geometry/UniverseAge.lean}.
\item[\phantomsection\label{lean:apparent-age}\textbf{HQIV apparent age.}]
  The companion snapshot $13.8\,\mathrm{Gyr}$ used as the
  wall-clock falsifier in \S\ref{sec:beta}. Lean:
  \leanid{age_apparent_Gyr_paper} in
  \leanfile{Hqiv/Geometry/UniverseAge.lean}.
\end{description}
\endgroup

\section{Reproducer script bundle}
\label{app:scripts}

The numerical results in this paper---the
$\beta_{\rm HQIV}=0.379^\circ$ prediction and its $-0.40\sigma$
agreement with Planck PR4 \cite{eskilt2022,planckPR42020}; the
$0.103^\circ$ apparent-age falsifier; the finite-mode Kirchhoff
energy density at FIRAS $T_{\rm CMB}$ \cite{fixsen2009}; and the
Friedmann $T^{2}$ recovery (Theorem~\ref{thm:friedmann})---are
reproduced by a small, dependency-light script bundle shipped
alongside the manuscript as \texttt{scripts.zip}. The bundle is
intended for upload as a Zenodo supplementary file together with
the paper PDF so that it shares the manuscript DOI and is
permanently citable.

\paragraph{Contents.}
\begingroup
\sloppy
\setlength{\emergencystretch}{4em}
\begin{description}
\item[\texttt{birefringence\_calculation.py}] Reproduces every line
  of the \S\ref{sec:beta} numerics paragraph: $m_T$,
  $\mathcal{N}(m_T)$, $t_{\rm wall}/t_{\rm Pl}$, $m_{\rm prop}$,
  $\beta_{\rm HQIV}$ in radians and degrees, residual to PR4 in
  $\sigma$ units, plus the apparent-age falsifier
  ($\beta^{\rm app}=0.103^\circ$, $-2.55\sigma$).
\item[\texttt{kirchhoff\_finite\_mode.py}] Computes the truncated
  blackbody energy density of \eqref{eq:Edens} on the FIRAS CMB
  temperature for a small UV/IR window, demonstrating the
  monotonicity statements of Theorem~\ref{thm:kirchhoff} and the
  hockey-stick cumulative count of \eqref{eq:cum} numerically.
\item[\texttt{friedmann\_recovery.py}] Confirms the
  Theorem~\ref{thm:friedmann} scaling
  $t_H^{\rm HQIV}\sim t_{\rm Pl}\,\mathcal{N}(m_T)\sim
  M_{\rm Pl}/T^{2}$ for several reference temperatures (CMB
  today, recombination, big-bang nucleosynthesis), and reports
  the relative discrepancy with the standard radiation-era
  Friedmann formula at each temperature.
\item[\texttt{README.md}] How-to-run instructions, contents
  table, patch-ontology disclaimer (the scripts are
  interoperability tools, not a continuum-limit refinement),
  and a citation block.
\end{description}
\endgroup

\paragraph{Requirements.}
Python 3.10+ with the standard library only.  No third-party
packages, no random seeds, and no external data files; the bundle
is fully deterministic and runs in well under a second on any
modern laptop.

\paragraph{Patch-ontology disclaimer.}
The numerical outputs are calculation-approximation translations
of patch quantities (cycle indices, shell ladders, normalised
readouts) into degrees, kelvin, and SI energies.  They are
\emph{not} a continuum-limit refinement of the patch theory; they
are the literature-translation rewrites called for by the
patch-theory reader contract above, packaged so that any reader
can re-execute them on their own machine.

\section*{Acknowledgments}
The author thanks the open-science community for the public
software stack on which the HQIV-LEAN library is built (Lean~4,
\texttt{mathlib}, and the Zenodo HQIV community
archive).\footnote{AI-assisted drafting and repository tooling
were used during manuscript preparation; all formal claims refer
to the HQIV-LEAN library, and the author reviewed hypothesis
lists in the cited Lean modules.}

\bibliographystyle{plain}
\bibliography{../references}

\end{document}
